% Save as: module6-problem-set.tex
% ============================================================================
% Problem Set 6: Causal Inference, RCTs & Structural + Experimental
% Integration
% Microeconomic Theory for Experimentalists & Applied Microeconomists
% Built from templates/problem_set_template.tex. Compile with pdflatex.
% ============================================================================
\documentclass[11pt]{article}
\usepackage{amsmath, amssymb, amsthm}
\usepackage[margin=1in]{geometry}
\usepackage{enumitem}
\usepackage{booktabs}

\newtheorem{question}{Question}

\title{Problem Set 6: Causal Inference, RCTs \& Structural + Experimental
Integration\\
\large Microeconomic Theory for Experimentalists \& Applied Microeconomists}
\author{Due: \textbf{[date --- before Lecture 1 of Module 7]}}
\date{}

\begin{document}
\maketitle

\noindent\emph{Ground rules: collaboration on ideas is encouraged; write-ups
are individual. You may use AI tools for parts flagged with
$\langle$AI$\rangle$ and must attach the transcript. Show all calculations.}

\medskip

\noindent\textbf{Weights:} Part A --- 30 points. Part B --- 45 points.
Part C --- 25 points.

\bigskip

% ---------------------------------------------------------------------------
\section*{Part A: Theory --- Guided Calculations (30 points)}

\begin{question}[Potential outcomes and selection bias, by hand --- 15
points]
Six workers face a voluntary training program. An omniscient observer
sees both potential outcomes (monthly earnings, in hundreds):

\begin{center}
\begin{tabular}{@{}lccc@{}}
\toprule
Person & $Y_0$ (untrained) & $Y_1$ (trained) & Enrolls voluntarily? \\
\midrule
1 & 30 & 50 & yes \\
2 & 50 & 60 & yes \\
3 & 70 & 70 & no  \\
4 & 40 & 60 & yes \\
5 & 60 & 70 & yes \\
6 & 80 & 80 & no  \\
\bottomrule
\end{tabular}
\end{center}

\begin{enumerate}[label=(\alph*)]
  \item \textbf{(3 pts)} Compute each person's treatment effect and the
        average treatment effect (ATE). Which people does the program
        help, and which does it leave unchanged?
  \item \textbf{(5 pts)} Under voluntary enrollment, the analyst
        observes $Y_1$ for enrollees and $Y_0$ for non-enrollees.
        Compute the naive comparison (mean observed outcome of enrollees
        minus mean observed outcome of non-enrollees) and note its sign.
  \item \textbf{(5 pts)} Decompose the naive comparison into the average
        effect on the treated (ATT) plus selection bias, computing each
        term numerically and verifying they sum to your answer in (b).
        \emph{Hint: selection bias compares the enrollees' and
        non-enrollees' UNTREATED outcomes.} One sentence: whose rational
        behavior created the bias?
  \item \textbf{(2 pts)} Now enrollment is assigned by coin flip. State
        what the difference in group means estimates, and name two
        questions about this program that even a perfectly executed RCT
        does NOT answer.
\end{enumerate}
\end{question}

\begin{question}[Power by simulation --- 15 points]
You are planning a correspondence study. Piloting suggests callback rates
of 9.7\% for the control name group and 6.5\% for the treatment name
group (the Bertrand--Mullainathan magnitudes). A classmate runs the Lab 6
power simulation (2{,}000 simulated experiments per sample size,
two-sided test at the 5\% level) and reports:

\begin{center}
\begin{tabular}{@{}lcccc@{}}
\toprule
$n$ per arm & 500 & 1{,}000 & 1{,}500 & 2{,}000 \\
Estimated power & 0.46 & 0.75 & 0.90 & 0.96 \\
\bottomrule
\end{tabular}
\end{center}

\begin{enumerate}[label=(\alph*)]
  \item \textbf{(4 pts)} In words a first-year student would understand:
        what does ``power $= 0.75$ at $n = 1{,}000$'' mean
        operationally, and why does power rise with $n$, rise with the
        true gap, and fall with outcome variance?
  \item \textbf{(4 pts)} Recommend a sample size for 80\% power,
        justifying your number from the table (interpolation is fine;
        state your reasoning). What total number of resumes does your
        recommendation imply, and why might you inflate it further
        before fielding the study?
  \item \textbf{(4 pts)} The budget caps the study at $n = 2{,}000$ per
        arm. The simulation reports that the standard error of the
        estimated gap at that size is about 0.0086. Compute the
        \emph{minimum detectable effect} at 80\% power (use
        $\text{MDE} \approx 2.8 \times \text{SE}$) and interpret it in
        one sentence: which true gaps would this study likely miss?
  \item \textbf{(3 pts)} Your advisor says: ``skip the simulation, use
        the standard formula.'' Give two concrete design features of a
        real correspondence study that the simulation can price and the
        textbook two-proportion formula cannot, and one sentence on why
        making assumptions explicit is a feature, not a cost.
\end{enumerate}
\end{question}

% ---------------------------------------------------------------------------
\section*{Part B: Experimental Design (45 points)}

\begin{question}[Separating taste-based from statistical discrimination]
Bertrand \& Mullainathan (2004) established the callback gap. Design a
correspondence study whose goal is the harder question: is the gap
driven by \emph{taste-based} discrimination (employers dislike hiring
the group, information cannot help) or \emph{statistical} discrimination
(group membership proxies for unobserved productivity, so individuating
information should make the proxy redundant)? Specify:
\begin{enumerate}[label=(\alph*)]
  \item \textbf{(6 pts)} The market: occupations, cities, how you find
        vacancies, and how many applications per vacancy you can send
        without detection --- and why your occupation choice matters for
        the taste-vs-statistical question.
  \item \textbf{(12 pts)} Treatments. Your design must cross the name
        manipulation with at least one \emph{information-richness}
        manipulation (e.g.\ verifiable certifications, references,
        licenses on the resume) and state the predicted callback
        pattern under each discrimination model --- including what the
        interaction term tells you that the two main effects cannot.
  \item \textbf{(9 pts)} Mechanics and ethics: how resumes are
        constructed and matched, how names are assigned, how callbacks
        are captured and coded, and the two main ethical costs of audit
        designs (employer deception and wasted screening time) with the
        standard mitigations an IRB will expect.
  \item \textbf{(9 pts)} Hypotheses and power: anchor on B\&M --- treat
        9.7\%/6.5\% as the \emph{baseline-arm} callback rates --- and
        your Part A(2) simulation logic; state the per-cell and total
        application counts your design implies once the information
        arms are added, and identify which comparison (main effect or
        interaction) drives the sample size.
  \item \textbf{(9 pts)} The biggest confound in YOUR design and your
        answer to it. (Candidates: names signal socioeconomic status,
        not only race; information manipulations change perceived
        overqualification; detection risk correlates with arm.)
\end{enumerate}
\end{question}

% ---------------------------------------------------------------------------
\section*{Part C: Silicon Pilot Interpretation $\langle$AI$\rangle$
(25 points)}

\begin{question}[Whose power calculation is it?]
You run \texttt{module6-simulation-lab.ipynb}. Your silicon hiring
managers produce a callback gap of 2.1 percentage points between name
groups, with very tight variance across runs. Your lab partner proposes:
``feed the silicon effect size and variance into our power calculator
--- it says we only need 800 resumes per arm, so we can cut the human
study's budget by half.''

\begin{enumerate}[label=(\alph*)]
  \item \textbf{(8 pts)} Give two distinct reasons the silicon effect
        size is not a valid stand-in for the human effect size (at
        least one specific to how LLMs are trained and aligned), and
        one reason the silicon \emph{variance} is even less trustworthy
        than the silicon mean.
  \item \textbf{(9 pts)} Write the power-calculation protocol you WOULD
        defend to an advisor: state where each input (effect size,
        base rate, variance/ICC, attrition) should come from, what role
        --- if any --- the silicon pilot legitimately plays in the
        protocol, and the sensitivity analysis you would attach.
  \item \textbf{(8 pts)} Now the twist: argue that your silicon result
        is not merely a failed pilot but a publishable object in its own
        right --- LLM resume screeners are deployed technology. State
        the research question your silicon audit actually answers, one
        robustness exercise it would need to be credible (the exercise
        may span multiple models and prompt framings), and the
        one-sentence scope limitation you would put in its abstract.
\end{enumerate}
Attach your AI transcripts and, if you ran code, the results CSV.
\end{question}

\end{document}
